\ExerciseSection

\begin{enumerate}
\def\labelenumi{\arabic{enumi})}
\item
  Let G be a discrete subgroup of rank p in \(R^{n}\) and let \((a_{i})_{1\leqslant i\leqslant p}\) be a free system of p points of G. The proof of Theorem 1 of no. 1 shows that G is a subgroup of the group generated by the p points \(d^{-1}a_{i}\) , where d is some integer. Show, without using Theorem 1, that there exists a free system of p points \(b_{i}=\sum_{j=1}^{p}b_{ij}a_{j}\) of G such that, if \(x_{i}=\sum_{j=1}^{p}x_{ij}a_{j}\quad(\mathrm{I}\leqslant i\leqslant p)\) is any free system of p points of G, we have \(|\det(x_{ij})|\geqslant|\det(b_{ij})|>0\) . Hence give a proof of Theorem 1 which does not rest on the theory of invariant factors, by showing that the \(b_{i}\) generate G. (Argue by contradiction: if \(z=\sum_{i=1}^{p}z_{i}b_{i}\in G\) and one of the \(z_{i}\) is not an integer, there exists a point \(u=\sum_{i=1}^{p}u_{i}b_{i}\) of the group generated by z and the \(b_{i}\) such that \(o<u_{i}<1\) for some index i, and hence obtain a contradiction.)
\item
  Let G be a discrete subgroup of \(R^{n}\) . If G is the direct sum of two subgroups H and K, then the intersection of the vector subspaces generated by H and K consists only of o, and the rank of G is therefore equal to the sum of the ranks of H and K (observe that, for every discrete subgroup G of \(R^{n}\) , the vector subspace of \(R^{n}\) generated by G is canonically isomorphic to \(G \otimes_{Z} R\) ).
\item
  Let G be a discrete subgroup of \(R^{n}\) . For a subgroup H of G to be a direct summand of G it is necessary and sufficient that \(H = V \cap G\) where V is a vector subspace of \(R^{n}\) (for necessity, use Exercise 2; for sufficiency, note that H is also the intersection of G and the vector subspace generated by H).
\item
  Let H and K be two discrete subgroups of \(R^{n}\) such that \(H + K\) is a closed subgroup. Show that \(H + K\) is then discrete and that \(r(\mathrm{H}) + r(\mathrm{K}) = r(\mathrm{H} \cap \mathrm{K}) + r(\mathrm{H} + \mathrm{K})\) . (Let V be the vector subspace of \(R^{n}\) generated by \(H \cap K\) ; split up H into the direct sum of \(V \cap H\) and a discrete group \(H_{1}\) , and K into the direct sum of \(V \cap K\) and a discrete group \(K_{1}\) , using Exercise 3; then show that the sum \(H_{1} + K_{1}\) is direct and use Exercise 2.)
\item
  Let G, \(G'\) be two closed subgroups of \(R^{n}\) such that \(G' \subset G\) , and let V (resp. \(V'\) ) be the largest vector subspace contained in G (resp. \(G'\) ). Show that there is a vector subspace W complementary to V such that G is the direct sum of V and the discrete group \(K = W \cap G\) , and such that \(G'\) is the direct sum of \(V \cap G'\) and \(K \cap G' = W \cap G'\) . (If U is a vector subspace complementary to \(V'\) , use Exercise 3 to show that the discrete group \(U \cap G'\) is the direct sum of \(U \cap V \cap G'\) and a discrete group \(K'\) , and take W to be a vector subspace containing \(K'\) .) Deduce that:
\end{enumerate}

\begin{enumerate}
\def\labelenumi{\alph{enumi})}
\tightlist
\item
  The quotient group \(G/G'\) is isomorphic to a product group of the form
\end{enumerate}

\[
\mathbf {R} ^ {p} \times \mathbf {T} ^ {q} \times \mathbf {Z} ^ {r} \times \mathbf {F},
\]

where F is a finite abelian group.

\begin{enumerate}
\def\labelenumi{\alph{enumi})}
\setcounter{enumi}{1}
\tightlist
\item
  Every closed subgroup and every Hausdorff quotient group of a group of the form \(R^{p} \times T^{q} \times Z^{r} \times F\) (F finite and abelian) are of the same form.
\end{enumerate}

\begin{enumerate}
\def\labelenumi{\arabic{enumi})}
\setcounter{enumi}{5}
\tightlist
\item
  Let H and K be two closed subgroups of \(R^{n}\) such that \(H + K\) is a closed subgroup.
\end{enumerate}

\begin{enumerate}
\def\labelenumi{\alph{enumi})}
\item
  Show that, if V (resp. W) is the largest vector subspace contained in H (resp. K), then \(V + W\) is the largest vector subspace contained in G, and therefore \(d(\mathrm{H}) + d(\mathrm{K}) = d(\mathrm{H} \cap \mathrm{K}) + d(\mathrm{H} + \mathrm{K})\) {[}note that the rational rank of \(\mathrm{G}/(\mathrm{V} + \mathrm{W})\) is finite and consequently that G cannot contain a line which is not contained in \(V + W\) {]}.
\item
  Let U be a subspace complementary to \(V + W\) , such that G is the direct sum of \(V + W\) and \(M = G \cap U\) , and such that \(H \cap K\) is the direct sum of \((V + W) \cap (H \cap K)\) and \(L = H \cap K \cap U\) (Exercise 5). Let \(H'\) (resp. \(K'\) ) be the subgroup of M consisting of the components of the points of H (resp. K) in the decomposition of G into the direct sum of \(V + W\) and M. Show that \(M = H' + K'\) and that \(H' \cap K' = L\) .
\item
  If we put \(H'' = H \cap (V + W)\) and \(K'' = K \cap (V + W)\) , show that \(r(H'') + r(K'') = r(H'' \cap K'') + r(H'' + K'')\) (reduce to the case where \(V \cap W = \{o\}\) , and show that then \(H'' \cap K''\) is the direct sum of \(V \cap K''\) and \(W \cap H''\) .
\end{enumerate}

\(d\) ) Show that

\[
r (\mathrm{H}) + r (\mathrm{K}) = r (\mathrm{H} \cap \mathrm{K}) + r (\mathrm{H} + \mathrm{K})
\]

{[}use c) and Exercise 4, noting that \(r(\mathbf{H}) = r(\mathbf{H}') + r(\mathbf{H}'')\) and \(r(\mathbf{K}) = r(\mathbf{K}') + r(\mathbf{K}'' )\) {]}.

\begin{enumerate}
\def\labelenumi{\arabic{enumi})}
\setcounter{enumi}{6}
\item
  Let G be a closed subgroup of \(R^{n}\) and let H be a closed subgroup of G. Then G is the direct sum of H and another closed subgroup K if and only if H is the intersection of G and a vector subspace. (Use Exercise 5 for the necessity; to show sufficiency, apply Theorem 2 of no. 2 appropriately to G and H, and use Exercise 3).
\item
  \begin{enumerate}
  \def\labelenumii{\alph{enumii})}
  \tightlist
  \item
    Let H, K be two closed subgroups of \(R^{n}\) . Show that if \(r(\mathrm{H}) + r(\mathrm{K})\) is equal to \(r(\mathrm{H} \cap \mathrm{K}) + r(\mathrm{H} + \mathrm{K})\) , then \(H + K\) is closed in \(R^{n}\) . {[}Using Exercise 7 and the hypothesis given, reduce to the case where H, K and \(H \cap K\) all generate the same vector subspace; if V (resp. W) denotes the largest subspace contained in H (resp. K), show with the help of Exercise 5 that V is generated by \(V \cap K\) , and W is generated by \(W \cap H\) ; finally, split up \(H \cap K\) into the direct sum of \(H \cap K \cap (V + W)\) and a discrete group, by using Exercise 7.{]}
  \end{enumerate}
\end{enumerate}

\begin{enumerate}
\def\labelenumi{\alph{enumi})}
\setcounter{enumi}{1}
\item
  Deduce from a) and Exercise 6 that, if H and K are two closed subgroups of \(R^{n}\) such that \(H + K\) is a closed subgroup, then \(H^{*} + K^{*}\) is also a closed subgroup.
\item
  Deduce from b) that if H and K are two closed subgroups of \(R^{n}\) such that \(d(\mathrm{H}) + d(\mathrm{K}) = d(\mathrm{H} \cap \mathrm{K}) + d(\overline{\mathrm{H} + \mathrm{K}})\) , then the subgroup \(H + K\) is closed.
\end{enumerate}

\begin{enumerate}
\def\labelenumi{\arabic{enumi})}
\setcounter{enumi}{8}
\item
  Let G be a subgroup of \(R^{n}\) , not necessarily closed, of rank p, and let V be the largest vector subspace contained in \(\overline{G}\) . If V has dimension \(q \leqslant p\) , show that G is the direct sum of \(V \cap G\) and a discrete subgroup of rank p - q contained in a vector subspace complementary to V {[}note that, for each \(x \in \overline{G}\) , \((x + V) \cap G\) is dense in the linear variety \(x + V\) {]}.
\item
  Let \(a\) be a real number and \(n\) an integer \(>0\) , and consider the numbers \(x_{k} = ka - [ka] \in [0, \mathrm{I}[(\mathrm{I} \leqslant k \leqslant n + \mathrm{I})\) . If \(\mathbf{I}_h\) denotes the interval
\end{enumerate}

\[
\left[ \frac {h - \mathrm{I}}{n}, \frac {h}{n} \right] \quad (\mathrm{I} \leqslant h \leqslant n),
\]

show that there exist two distinct indices \(k, k'\) such that \(x_k\) and \(x_{k'}\) belong to the same interval \(I_h\) (for a suitable value of \(h\) ). Deduce that there exist two integers \(p, q\) such that \(1 \leqslant p \leqslant n\) and \(|pa - q| \leqslant 1 / n\) .

\begin{enumerate}
\def\labelenumi{\Roman{enumi})}
\setcounter{enumi}{1}
\tightlist
\item
  Let \(\boldsymbol{a}_i = (a_{ij})\) ( \(1 \leqslant i \leqslant m\) , \(1 \leqslant j \leqslant n\) ) be \(m\) points of \(\mathbf{R}^n\) and let \(q\) be an integer \(>0\) . For each point \(\boldsymbol{k} = (k_i)\) ( \(1 \leqslant i \leqslant m\) ) of \(\mathbf{R}^m\) with integer coordinates, not all zero, satisfying the inequalities \(0 \leqslant k_i \leqslant q\) , consider the point \(x_k = (x_{j,k})\) ( \(1 \leqslant j \leqslant n\) ) of \(\mathbf{R}^n\) , all of whose coordinates belong to {[}o, i{[} and which is congruent mod \(Z^n\) to \(\sum_{i=1}^{m} k_i a_i\) , i.e., the point with coordinates
\end{enumerate}

\[
x _ {j, k} = \sum_ {i = 1} ^ {m} k _ {i} a _ {i j} - \left[ \sum_ {i = 1} ^ {m} k _ {i} a _ {i j} \right].
\]

If \(p\) is the smallest integer such that \(q + 1 \geqslant p^{n/m}\) , show that there are two distinct points \(k, k'\) such that \(x_k\) and \(x_{k'}\) both belong to the same cube of side \(1/p\) (same method as Exercise 10). Deduce that there exist \(m\) integers \(p_i\) , not all zero, and \(n\) integers \(r_j\) ( \(1 \leqslant j \leqslant n\) ) such that \(0 \leqslant p_i \leqslant q\) for \(1 \leqslant i \leqslant m\) , and

\[
\left| \sum_ {i = 1} ^ {m} p _ {i} a _ {i j} - r _ {j} \right| \leqslant \frac {\mathrm{I}}{p} \quad (\mathrm{I} \leqslant j \leqslant n).
\]

Use this result to give another proof of Proposition 2 of no. 1 (the ``pigeon-hole principle'').

\begin{enumerate}
\def\labelenumi{\arabic{enumi})}
\setcounter{enumi}{11}
\tightlist
\item
  For each pair of real numbers \((\theta, \beta)\) there exists an infinite number of triples \((p, q, r)\) of integers such that \(r > 0\) , \(|q| \leqslant \frac{1}{2} r\) and
\end{enumerate}

\[
- \frac {\mathrm{I}}{r} <   \theta q + p - \beta <   \frac {\mathrm{I}}{r}.
\]

{[}If \(m, n\) are two integers such that \(|n\theta - m| < 1/n\) (Exercise 10), take \(r = n\) and choose \(p, q\) so that \(pn + qm\) differs from \(n\beta\) by less than \(1/2\) {]}.

\begin{enumerate}
\def\labelenumi{\arabic{enumi})}
\setcounter{enumi}{12}
\tightlist
\item
  The Farey series of order \(n\) ( \(n\) an integer \(>0\) ) is the set \(F_n\) of rational numbers \(p/q\) in their lowest terms such that \(0 \leqslant p \leqslant q \leqslant n\) , arranged in increasing order (i.e., the Farey series is a sequence).
\end{enumerate}

\begin{enumerate}
\def\labelenumi{\alph{enumi})}
\item
  Show that if two rational numbers \(r = p / q\) , \(r' = p' / q'\) are such that \(p'q - pq' = \pm 1\) , then every pair of integers \((p'', q'')\) can be expressed in the form \(p'' = px + p'y\) , \(q'' = qx + q'y\) where \(x\) and \(y\) are integers. The fraction \(p'' / q''\) belongs to the closed interval with end-points \(r\) and \(r'\) if and only if \(x\) and \(y\) have the same sign.
\item
  Deduce from a) that if \(r = p / q\) and \(r' = p' / q'\) are two rational numbers in {[}0, 1{]} such that \(q > 0\) and \(q' > 0\) and \(qp' - pq' = \pm 1\) , then \(r\) and \(r'\) are consecutive elements in the sequence \(\mathbf{F}_n\) , where \(n = \max (q, q')\) . Furthermore, the smallest integer \(m\) such that the open interval with end-points \(r\) and \(r'\) contains a point of \(\mathbf{F}_m\) is \(m = q + q'\) , and there is only one point of \(\mathbf{F}_{q + q'}\) in this interval, namely the fraction \((p + p') / (q + q')\) .
\item
  Show that conversely, if \(r\) and \(r'\) are two consecutive elements of \(\mathbf{F}_n\) , we have \(p'q - pq' = \pm 1\) (use induction on \(n\) ).
\item
  Deduce from c) that, for each real number \(\theta \in [\mathrm{o},\mathrm{i}]\) and each integer \(n\geqslant \mathrm{i}\) , there is at least one fraction \(p / q\) in its lowest terms such that \(\mathrm{i}\leqslant q\leqslant n\) and \(|\theta -p / q|\leqslant \mathrm{i} / (n + \mathrm{i})q\) (cf.~Exercise 10).
\item
  If \(p / q\) is a fraction in its lowest terms such that \(|\theta - p / q| < 1 / q^2\) , then \(\theta\) belongs to the open interval whose end-points are the two terms of the Farey series \(\mathbf{F}_q\) consecutive to \(p / q\) .
\end{enumerate}

\begin{enumerate}
\def\labelenumi{\arabic{enumi})}
\setcounter{enumi}{13}
\tightlist
\item
  Let \(\varphi: \mathbf{R} \to \mathbf{T}\) be the canonical homomorphism; let \(\theta\) be an element of infinite order in \(\mathbf{T}\) and let \(\theta_0\) be an (irrational) real number such that \(\varphi(\theta_0) = \theta\) . For each integer \(n > 0\) , let \(\mathbf{S}_n\) be the set of elements \(k\theta\) of \(\mathbf{T}\) for which \(\mathbf{i} \leqslant k \leqslant n\) , and for each interval \(\mathbf{I}\) of \(\mathbf{R}\) let \(\mathrm{N}(\mathbf{I}, n)\) be the number of elements of the set \(\mathbf{I} \cap \overline{\varphi}^1(\mathbf{S}_n)\) .
\end{enumerate}

\begin{enumerate}
\def\labelenumi{\alph{enumi})}
\item
  If we take \(\mathbf{I} = [\mathrm{o},\theta_0[\) , show that \(\mathrm{N(I,n)} = [n\theta_0]\) {[}note that \(\mathrm{N(I,n)}\) is then the number of pairs of integers \((x,y)\) such that \(\mathrm{i}\leqslant y\leqslant n\) and \(x\leqslant y\theta_0 < x + \theta_0]\) .
\item
  More generally, if we take \(\mathbf{I} = [\mathrm{o}, m\theta_0[, m\) being an integer, we have \(m. [(n - m)\theta_0] \leqslant \mathrm{N}(\mathrm{I}, n) \leqslant m. [n\theta_0]\) (same method).
\item
  Deduce that, if I is any interval, the number \(\mathrm{N}(\mathrm{I}, n)/n\) tends to a limit equal to the length of I, as \(n \to +\infty\) . (Prove the result first when I is of the form \([m\theta_{0} + a, m'\theta_{0} + a']\) , where \(m, m', a\) and \(a'\) are integers, by using b); then pass to the general case by approximating to the endpoints of I by numbers of the form \(m\theta_{0} + a\) .) {[}``Equipartition of the sequence \((k\theta) \mod 1\)''.{]}
\end{enumerate}

* 15) Let I be a closed cube in \(\mathbf{R}^n\) with side \(2\pi\) . Map each point \(x = (x_i)\) of I to the point \(y = (y_j)\) of \(\mathbf{R}^{n+1}\) such that

\[
\begin{array}{l} y _ {1} = \sin x _ {1} \\ y _ {2} = (2 + \cos x _ {1}) \sin x _ {2} \end{array}
\]

\[
y _ {p} = \left(2 ^ {p - 1} + \sum_ {k = 1} ^ {p - 1} 2 ^ {p - k - 1} \cos x _ {p - k} \cos x _ {p - k + 1} \dots \cos x _ {p - 1}\right) \sin x _ {p} \quad (2 \leqslant p \leqslant n)
\]

\[
y _ {n + 1} = \sum_ {k = 1} ^ {n - 1} 2 ^ {p - k - 1} \cos x _ {n - k} \cos x _ {n - k + 1} \dots \cos x _ {n}.
\]

Show that the image of I under this mapping is homeomorphic to \(T^{n}\) (proof by induction on n: observe that \(y_{p}\) has the sign of \(\sin x_{p}\) for \(2 \leqslant p \leqslant n\)). If n=2, the subset of \(R^{3}\) so defined is called a torus of revolution.*

¶ 16) Let G be a subgroup of \(\mathbf{R}^n\) which contains a compact connected subset K such that the affine linear variety generated by K is of dimension \(p\) . Show that G contains a vector subspace of dimension \(p\) . {[}Proof by induction on \(n\) , using Exercise 10 c) of Chapter VI, § 1.{]}

\(\mathbb{J}\) 17) Let \(\mathbf{E}\) be the product \((\mathbf{Q}_p)^n\) of \(n\) factors equal to the field \(\mathbf{Q}_p\) of \(p\) -adic numbers (Chapter III, § 6, Exercise 23), endowed with the topological group structure which is the product of the additive topological structures of the factors \(\mathbf{Q}_p\) , and with its \(n\) -dimensional vector space structure over the field \(\mathbf{Q}_p\) . If \(v\) denotes the additive \(p\) -adic valuation on \(\mathbf{Q}_p\) , put \(|x|_p = p^{-v_p(x)}\) for all \(x \neq 0\) in \(\mathbf{Q}_p\) , and \(|o|_p = 0\) ; and put \(||x|| = \sup |x_i|_p\) for all \(x = (x_i)_{1 \leqslant i \leqslant n} \in \mathbf{E}\) (cf.~Chapter IX, § 3, nos. 2 and 3).

\begin{enumerate}
\def\labelenumi{\alph{enumi})}
\tightlist
\item
  A subset \(\mathbf{A}\) of \(\mathbf{E}\) is relatively compact if and only if
\end{enumerate}

\[
\sup _ {\mathbf {x} \in \mathbf {A}} | | \mathbf {x} | | <   + \infty .
\]

{[}Note that \(v_{p}\) is a continuous mapping of \(\mathbf{Q}_p\) into \(\mathbf{Z}\) (with the discrete topology).{]}

\begin{enumerate}
\def\labelenumi{\alph{enumi})}
\setcounter{enumi}{1}
\item
  If \(\mathbf{G}\) is a closed subgroup of \(\mathbf{E}\) , then \(\mathbf{G}\) is a topological module over the ring \(\mathbf{Z}_p\) of \(p\) -adic integers (if \(\boldsymbol{x} \in \mathbf{G}\) , we have \(n\boldsymbol{x} \in \mathbf{G}\) for all \(n \in \mathbf{Z}\) , and \(\mathbf{Z}\) is dense in \(\mathbf{Z}_p\) ).
\item
  If K is a compact subgroup of E, there exists a free system of \(m \leqslant n\) points \(a_{i}\) ( \(i \leqslant i \leqslant m\) ) of E, such that K is the direct sum of the m groups \(Z_{p}.a_{i}\) . {[}Use a) to show that, if \((\boldsymbol{e}_{j})_{1 \leqslant j \leqslant n}\) is the canonical base of E, there is an integer \(k \in Z\) such that K is contained in the direct sum of the n groups \(Z_{p}.p^{k}\boldsymbol{e}_{j}\) ; then apply the theory of modules over a principal ideal ring.{]}
\item
  If G is a closed (but not compact) subgroup of E, then G contains a vector subspace of dimension 1, \(Q_{p} \cdot a\) with \(a \neq 0\) . {[}Let C be the compact subset of E consisting of all points \(x \in E\) such that \(\|x\| = 1\) ; show that \(G \cap C\) contains a sequence of points \((\boldsymbol{x}_{r})_{r \in \mathbb{N}}\) such that \(p^{-r}x_{r} \in G\) , and take a to be a cluster point of this sequence.{]}
\item
  Let G be a closed subgroup of E, let V be the largest vector subspace contained in G and let W be a vector subspace complementary to V; show that \(W \cap G\) is compact and that G is the direct sum of V and \(W \cap G\) (argue as in Theorem 2 of no. 2).
\end{enumerate}

\(f)\) Given a subgroup \(\mathbf{G}\) of \(\mathbf{E}\) , let \(\mathbf{G}^*\) denote the set of all \(\boldsymbol{u} = (u_i) \in \mathbf{E}\) such that

\[
(u | x) = \sum_ {i = 1} ^ {n} u _ {i} x _ {i}
\]

belongs to \(\mathbf{Z}_p\) for all \(x = (x_i) \in \mathbf{G}\) . Show that \(\mathbf{G}^*\) is closed in \(\mathbf{E}\) , that \((\overline{\mathbf{G}})^* = \mathbf{G}^*\) and that \((\mathbf{G}^*)^* = \overline{\mathbf{G}}\) .

\begin{enumerate}
\def\labelenumi{\alph{enumi})}
\setcounter{enumi}{6}
\tightlist
\item
  State and prove the analogues of Exercises 2 to 8 inclusive for closed subgroups of E. In particular, every closed subgroup and every Hausdorff quotient group of a group of the form
\end{enumerate}

\[
(\mathbf {Q} _ {p}) ^ {r} \times (\mathbf {Q} _ {p} / \mathbf {Z} _ {p}) ^ {s} \times (\mathbf {Z} _ {p}) ^ {t} \times \mathbf {F}
\]

(where F is the product of a finite number of cyclic groups of p-power order) is a group of the same form.

